\chapter{Limitações e trabalho futuro}
\label{cap:limitacoes}

O desenvolvimento do presente projeto permitiu construir uma base funcional
para a representação e utilização de conhecimento para o apoio à tomada de
decisão em veículos autónomos. No entanto, no momento atual e tendo em conta
o estado de desenvolvimento do projeto, existem limitações que devem ser
consideradas na análise dos resultados obtidos, até à data. Estas limitações
estão relacionadas tanto com o processo de aquisição do conhecimento como com
as funcionalidades que ainda não se encontram totalmente implementadas.

\section{Limitações}
\label{sec:limitacoes}

Uma das principais limitações foi o momento em que foi possível ter o
primeiro contacto com o perito. Uma parte significativa do trabalho já se
encontrava em desenvolvimento quando foram realizadas as primeiras reuniões.
Até esse momento, as decisões relativas à representação do ambiente, aos
factos, às regras e aos possíveis cenários tinham sido definidas
principalmente com base na pesquisa realizada e no conhecimento desenvolvido
de cada elemento do grupo. Neste sentido, algumas opções foram, inicialmente,
decididas em grupo, sem existir qualquer validação especializada para
confirmar se o caminho que o grupo estava a seguir era o mais adequado. O
contacto posterior com o perito permitiu esclarecer algumas destas questões e
identificar aspetos que deveriam ser revistos, nomeadamente a importância da
relação entre velocidade, distância e capacidade de travagem, o número de
leituras a considerar, as prioridades nas decisões e a permanência dos factos
na memória de trabalho. Contudo, como apenas se realizaram duas reuniões com
o perito, algumas decisões atualmente presentes no projeto são ainda soluções
em evolução, passíveis de análise e validação futuras.

Outra limitação está relacionada com a própria evolução da base de
conhecimento. À medida que são introduzidas novas situações, exceções e
relações entre diferentes elementos do domínio, aumenta também o número de
regras que é necessário manter e articular. Embora a organização dos
ficheiros DRL, por áreas de conhecimento, facilite este processo, com o
crescimento da base de regras pode ser mais difícil identificar sobreposições,
conflitos ou condições não previstas. A revisão sistemática das regras, o
recurso aos cenários de teste e a continuação da validação com o perito serão,
por isso, importantes à medida que o sistema evoluir.

Relativamente à própria natureza da simulação, existem também algumas
limitações. O objetivo do sistema não é reproduzir de forma completa a
dinâmica de um veículo autónomo real. O ambiente, a movimentação do veículo,
as vias de circulação e os restantes elementos rodoviários estão
representados de uma maneira simplificada, uma vez que o foco do trabalho é o
desenvolvimento do motor de inferência e da base de conhecimento. Neste
sentido, questões como o comportamento real dos pneus ou outras condições mais
complexas não serão representadas com o mesmo nível de detalhe que seria
necessário, caso o sistema fosse para utilização real. Também os sensores são
simulados e recebem informação estruturada a partir do ambiente de simulação.
Desta forma, o projeto não aborda problemas associados à perceção real como,
por exemplo, o processamento de imagem.

A explicação das decisões carece de complexidade. O \emph{Why} e o \emph{How}
estão mais desenvolvidos, uma vez que o sistema consegue associar uma decisão
à regra que lhe deu origem, à razão apresentada e à evidência utilizada, bem
como acompanhar alterações na memória de trabalho e a execução das regras. O
\emph{Why Not}, pelo contrário, necessita ainda de ser aperfeiçoado. Uma vez
que o objetivo é explicar por que motivo uma determinada alternativa não foi
escolhida, é necessário conservar as propostas concorrentes que não foram
selecionadas e, quando aplicável, identificar as condições que impediram
determinada regra de resultar numa proposta.

A componente em Prolog é outra limitação do estado atual do trabalho, por
enquanto, uma vez que foi dada prioridade inicial à implementação em Drools.
Apesar de esta decisão permitir consolidar primeiro a arquitetura, os factos,
os eventos e as regras, a comparação entre os dois paradigmas ainda não se
encontra definida.

\section{Trabalho futuro}
\label{sec:trabalho-futuro}

Uma das prioridades futuras será continuar a revisão da base de conhecimento
com o apoio do perito. Uma vez que o primeiro contacto se deu numa fase já
relativamente avançada do desenvolvimento, será necessário verificar se as
condições, prioridades e ações representam, adequadamente, o conhecimento do
domínio que temos vindo a ter em conta. Neste contexto, um dos aspetos
relevantes a aprofundar é a forma como a criticidade é determinada. A
integração explícita da velocidade, da distância ao objeto, da velocidade
relativa, do tempo de reação e da capacidade de travagem poderá permitir
substituir alguns limiares fixos por critérios mais dependentes do contexto
observado. Ao nível da informação sensorial, uma evolução importante consiste
na introdução de níveis de confiança associados às deteções e na definição de
regras de confirmação entre diferentes fontes. Desta forma, o sistema poderá
distinguir não só o valor de uma determinada observação, como também o grau
de confiança da mesma.

A componente que dá resposta ao \emph{Why Not} necessita de ser desenvolvida
com mais pormenor e atenção. Para cada alternativa não escolhida poderá ser
registado o respetivo motivo, como uma prioridade inferior ou uma condição não
satisfeita. Esta informação permitirá apresentar explicações construtivas,
mais claras e mais úteis para a análise e validação da base de conhecimento.

Para além da implementação em Prolog, podem ainda ser introduzidas situações
com maior interação entre veículos, obstáculos inesperados, diferentes
condições meteorológicas e combinações mais complexas de sinalização, com o
objetivo de acrescer complexidade à simulação apresentada. A Teoria dos Jogos
poderá também ser explorada como abordagem complementar em situações onde a
decisão depende do comportamento esperado de outros intervenientes,
nomeadamente nas mudanças de via. Neste caso, as regras de segurança do
sistema pericial continuariam a funcionar como restrições, enquanto uma
componente adicional seria responsável por apoiar a escolha entre as
diferentes manobras consideradas seguras e possíveis. Posteriormente, poderá
ainda ser considerado um ambiente de simulação mais detalhado ou a utilização
de dados provenientes de sensores reais. No entanto, essa evolução deverá
preservar o objetivo principal do projeto.

Assim, as limitações identificadas não dizem respeito apenas às
funcionalidades ainda incompletas, mas também ajudam a definir as etapas
seguintes do desenvolvimento. A continuação do contacto com o perito, a
revisão das regras, a conclusão do \emph{Why Not} e do Prolog e a introdução
de novos cenários permitirão aproximar progressivamente o protótipo do modelo
conceptual.
